\documentclass[12pt,table,xcolor={dvipsnames}]{beamer}
\usetheme{Pittsburgh}
\usecolortheme{seagull}
%\usepackage[utf8]{inputenc}
\usepackage{fontspec}
\usepackage[brazilian]{babel}
\usepackage{amsmath}
\usepackage{listings}
\usepackage{amsfonts}
\usepackage{amssymb}
\usepackage{graphicx}
\author{Introdução ao C++}
\title{Estruturas de Dados}
%\setbeamercovered{transparent}
\setbeamertemplate{navigation symbols}{}
%\logo{\includegraphics[scale=0.015]{Brasao_UFSC.png}\includegraphics[scale=0.2]{brasao_PPGCC.jpg}}
\institute{Departamento de Informática e de Estatística \\ Prof. Jean Everson Martina \\ Prof. Aldo von Wangenheim}
\date{2016.2}
\subject{}
\usebackgroundtemplate{\includegraphics[width=\paperwidth,
height=\paperheight]{../reusable_images/fundo_UFSC.png}}
\begin{document}

{
\usebackgroundtemplate{\includegraphics[width=\paperwidth,
height=\paperheight]{../reusable_images/fundo_capa.png}}
\begin{frame}
\titlepage
\includegraphics[scale=0.3]{../reusable_images/brasao_INE.png}
\end{frame}
}

\begin{frame}{Visão Geral}
\begin{itemize}
\item Histórico da Linguagem;
\item Programa C++: header, source – função main();
\item GCC/G++;
\item Arquivos objeto, bibliotecas dinâmicas e estáticas;
\item \#include, \#define, namespaces, typedefs;
\item Ponteiros, Referências, arrays, primitivas, estruturas de controle;
\end{itemize}
\end{frame}

\begin{frame}{Características do C++}
\begin{itemize}
\item Linguagem multi-paradigma;
\item Programação Procedural;
\item Orientação a objetos;
\item Meta Programação;
\item Desalocação memória – Decisão cabe ao programador;
\end{itemize}
\end{frame}


\begin{frame}{Histórico}
\begin{itemize}
\item Criada por Bjarne Stroustrup em 1979 no Bell Labs da AT\&T e inicialmente chamada de ``C with classes”;
\item Em 1983 passou a ser chamada C++;
\item Ao longo do tempo incorporou novas funcionalidades como:
\begin{itemize}
\item Herança múltipla;
\item Sobreescrita de operadores;
\item Templates;
\item Entre outras coisas que veremos ao longo do curso.
\end{itemize}
\end{itemize}
\end{frame}

\begin{frame}{Histórico}
\begin{itemize}
\item Possui compatibilidade com C;
\item Um compilador C++ deve compilar um programa escrito em C;
\item Existem vários comiladores da linguagem;
\item É um linguagem padronizada internacionamelmente:
\begin{itemize}
\item Em 1998 foi definido um padrão ISO para a linguagem;
\item Revisado em 2003;
\item Revisado em 2011;
\item Revisado em 2014;
\end{itemize}
\end{itemize}
\end{frame}

\begin{frame}{GCC/G++}
\begin{itemize}
\item GNU Compiler Collection;
\item Compilador originalmente desenvolvido focando a linguagem C, chamava-se GNU C
Compiler;
\item Ao longo do anos passou a contemplar outras linguagens, como C++;
\end{itemize}
\end{frame}

\begin{frame}{Programa C++ - Arquivos}
\begin{itemize}
\item Header - Definições visíveis a outros headers e ao compilador;
\item Sources - Corpo de métodos e funções;
\item Implemented Headers - Headers com implementado que tem que ser instrumentada pelo compilador;
\item Função Main - Função inicial invocada pelo executável;
\end{itemize}
\end{frame}

\begin{frame}[fragile]{Header Exemplo}
\lstset{language=C++,
          keywordstyle=\color{blue}\ttfamily,
          stringstyle=\color{red}\ttfamily,
          commentstyle=\color{green}\ttfamily,
          breaklines=true,
          basicstyle=\ttfamily\footnotesize
          }
\begin{lstlisting}
// in myclass.h
#ifndef MYCLASS_H
#define MYCLASS_H

class MyClass {
 public:
  void foo();
  float bar;
};
#endif
\end{lstlisting}
\end{frame}

\begin{frame}[fragile]{Source Exemplo}
\lstset{language=C++,
          keywordstyle=\color{blue}\ttfamily,
          stringstyle=\color{red}\ttfamily,
          commentstyle=\color{green}\ttfamily,
          breaklines=true,
          basicstyle=\ttfamily\footnotesize
          }
\begin{lstlisting}
// in myclass.cpp
#include "myclass.h"
#include <iostream>

void MyClass::foo() {
  std::cout<<"Hello World!"<<std::endl;
}
\end{lstlisting}
\end{frame}


\begin{frame}[fragile]{HPP Exemplo}
\lstset{language=C++,
          keywordstyle=\color{blue}\ttfamily,
          stringstyle=\color{red}\ttfamily,
          commentstyle=\color{green}\ttfamily,
          breaklines=true,
          basicstyle=\ttfamily\footnotesize
          }
\begin{lstlisting}
// in myclass.hpp
#include <iostream>

template <class T>
class MyClass {
 public:
  void foo(){
   std::cout<<"Hello World!"<<std::endl;
  }
  T bar;
};
\end{lstlisting}
\end{frame}

\begin{frame}[fragile]{Programa Exemplo}
\lstset{language=C++,
          keywordstyle=\color{blue}\ttfamily,
          stringstyle=\color{red}\ttfamily,
          commentstyle=\color{green}\ttfamily,
          breaklines=true,
          basicstyle=\ttfamily\footnotesize
          }
\begin{lstlisting}
//in main.cpp
//#include "myclass.h"  // defines MyClass
#include "myclass.hpp"  // defines MyClass

int main()
{
  //MyClass* a = new MyClass();
  MyClass<int>* a = new MyClass<int>();
  a->foo();
  return 0;
}
\end{lstlisting}
\end{frame}

\begin{frame}[fragile]{Compilando Programa Exemplo}
\lstset{language=C++,
          keywordstyle=\color{blue}\ttfamily,
          stringstyle=\color{red}\ttfamily,
          commentstyle=\color{green}\ttfamily,
          breaklines=true,
          basicstyle=\ttfamily\footnotesize
          }
\begin{lstlisting}
$ g++ main.cpp -o hello_world -std=c++11 -Werror
$ ./hello_world
$ Hello World!
\end{lstlisting}
\end{frame}


\begin{frame}[fragile]{Diretiva \#include}
\begin{itemize}
\item Informações sobre arquivos, símbolos e tipos interpretadas no pré-processamento da compilação;
\item Avisa ao compilador quais headers ele precisa carregar para pré-processar as instruções do arquivo em questão, mais ou menos como o import do Java;
\item Pode ser utilizado tanto no header quanto no source;
\item Evite \#includes desnecessários no header, isso torna a etapa de análise mais custosa.
\end{itemize}
\end{frame}

\begin{frame}[fragile]{Diretiva \#include}
\lstset{language=C++,
          keywordstyle=\color{blue}\ttfamily,
          stringstyle=\color{red}\ttfamily,
          commentstyle=\color{green}\ttfamily,
          breaklines=true,
          basicstyle=\ttfamily\footnotesize
          }
\begin{itemize}
\item \#include pode ser usar para escolher código durante a compilação;
\begin{lstlisting}
#include "structs/my_struct.h"
 int main(){
  myStruct ms;
 }
\end{lstlisting}
\item \$ g++ include.cc -std=c++11 -Werror;
\begin{lstlisting}
#include "my_struct.h"
 int main(){
  myStruct ms;
 }
\end{lstlisting}
\item \$ g++ include.cc -I structs -std=c++11 -Werror.
\end{itemize}
\end{frame}

\begin{frame}[fragile]{Diretiva \#include}
\lstset{language=C++,
          keywordstyle=\color{blue}\ttfamily,
          stringstyle=\color{red}\ttfamily,
          commentstyle=\color{green}\ttfamily,
          breaklines=true,
          basicstyle=\ttfamily\footnotesize
          }
\begin{itemize}
\begin{lstlisting}
#include <list>
\end{lstlisting}
\item Procura o header nos diretórios do sistema e aqueles passados como parâmetro ao compilador.
\\
\begin{lstlisting}
#include "my_struct.h"
\end{lstlisting}
\item Procura o header primeiro no diretório que contém o arquivo sendo compilado, e depois nos diretórios do sistema e aqueles passados como parâmetro ao compilador.
\end{itemize}
\end{frame}

\begin{frame}[fragile]{Diretiva \#define}
\lstset{language=C++,
          keywordstyle=\color{blue}\ttfamily,
          stringstyle=\color{red}\ttfamily,
          commentstyle=\color{green}\ttfamily,
          breaklines=true,
          basicstyle=\ttfamily\footnotesize
          }
\begin{itemize}
\item Atribui um valor a um símbolo;
\item Cada vez que o pré-processador encontra este símbolo no código ele substitui o símbolo por esse valor.
\begin{lstlisting}
#define max_iterations 1000
\end{lstlisting}
\end{itemize}
\end{frame}

\begin{frame}[fragile]{Diretiva \#define}
\lstset{language=C++,
          keywordstyle=\color{blue}\ttfamily,
          stringstyle=\color{red}\ttfamily,
          commentstyle=\color{green}\ttfamily,
          breaklines=true,
          basicstyle=\ttfamily\footnotesize
          }
\begin{itemize}
\item Também pode ser usado para definir "macros", pequenos trechos de código que são colados cada vez que o pré-processador identifica o símbolo associado.
\begin{lstlisting}
#define getmax(a,b) ((a)>(b)?(a):(b))

int main() {
 int x = 10;
 int y = getmax(x,5); //y vale 10
}
\end{lstlisting}
\end{itemize}
\end{frame}

\begin{frame}[fragile]{Diretiva \#define}
\lstset{language=C++,
          keywordstyle=\color{blue}\ttfamily,
          stringstyle=\color{red}\ttfamily,
          commentstyle=\color{green}\ttfamily,
          breaklines=true,
          basicstyle=\ttfamily\footnotesize
          }
\begin{itemize}
\item Com a sustituição feita ele fica assim:
\begin{lstlisting}
#define getmax(a,b) ((a)>(b)?(a):(b))

int main() {
 int x = 10;
 int y =  ((x)>(5)?(x):(5));; //y vale 10
}
\end{lstlisting}
\end{itemize}
\end{frame}

\begin{frame}[fragile]{Diretiva \#define}
\lstset{language=C++,
          keywordstyle=\color{blue}\ttfamily,
          stringstyle=\color{red}\ttfamily,
          commentstyle=\color{green}\ttfamily,
          breaklines=true,
          basicstyle=\ttfamily\footnotesize
          }
\begin{itemize}
\item So devemos usar quando o código é pequeno e existe um real conhecimento do fluxo das instruções;
\item Desgraças podem acontecer e o código colado pelo pré-processador pode tomar um rumo inesperado;
\item Como o código é "colado", não existe ônus de ter que colocar o endereço da instrução na pilha como em uma chamada de função;
\item Em uma próxima aula discutiremos isso ao ver funções inline.
\end{itemize}
\end{frame}

\begin{frame}[fragile]{Diretiva \textit{namespace}}
\lstset{language=C++,
          keywordstyle=\color{blue}\ttfamily,
          stringstyle=\color{red}\ttfamily,
          commentstyle=\color{green}\ttfamily,
          breaklines=true,
          basicstyle=\ttfamily\footnotesize
          }
\begin{itemize}
\item Estabelecem o domínio ao qual declarações (classes, structs, metaclasses) fazem parte;
\item É utilizado para organizar o código em diferentes domínios, que dependendo da
estrutura de diretórios do código fonte;
\item Apresentam uma leve semelhança com a organização em pacotes do java;
\item Também serve para eliminar ambigüidades.
\end{itemize}
\end{frame}

\begin{frame}[fragile]{Diretiva \textit{namespace}}
\lstset{language=C++,
          keywordstyle=\color{blue}\ttfamily,
          stringstyle=\color{red}\ttfamily,
          commentstyle=\color{green}\ttfamily,
          breaklines=true,
          basicstyle=\ttfamily\footnotesize
          }
\begin{itemize}
\begin{lstlisting}
#include "myclass.h"
#include <iostream>
void MyClass::foo() {
  std::cout<<"Hello World!"<<std::endl;
}\end{lstlisting}
\item Passamos a não precisar chamar as funcões com seu nome completo
\begin{lstlisting}
#include "myclass.h"
#include <iostream>

using namespace std;
void MyClass::foo() {
  cout<<"Hello World!"<<endl;
}
\end{lstlisting}

\end{itemize}
\end{frame}

\begin{frame}[fragile]{HPP Exemplo \textit{namespace}}
\lstset{language=C++,
          keywordstyle=\color{blue}\ttfamily,
          stringstyle=\color{red}\ttfamily,
          commentstyle=\color{green}\ttfamily,
          breaklines=true,
          basicstyle=\ttfamily\footnotesize
          }
\begin{lstlisting}
// in myclass.hpp
#include <iostream>

using namespace templates;

template <class T>
class MyClass {
 public:
  void foo(){
   std::cout<<"Hello World!"<<std::endl;
  }
  T bar;
};
\end{lstlisting}
\end{frame}

\begin{frame}[fragile]{Programa Exemplo \textit{namespace}}
\lstset{language=C++,
          keywordstyle=\color{blue}\ttfamily,
          stringstyle=\color{red}\ttfamily,
          commentstyle=\color{green}\ttfamily,
          breaklines=true,
          basicstyle=\ttfamily\footnotesize
          }
\begin{lstlisting}
//in main.cpp
#include "myclass.h"  // defines MyClass
#include "myclass.hpp"  // defines MyClass

using namespace std;

int main()
{
  MyClass* a; //ambiguidade
  MyClass<int>* b; //ambiguidade
}
\end{lstlisting}
\begin{itemize}
\item Qual das implementacões de MyClass foi usada?
\end{itemize}
\end{frame}

\begin{frame}[fragile]{Programa Exemplo \textit{namespace}}
\lstset{language=C++,
          keywordstyle=\color{blue}\ttfamily,
          stringstyle=\color{red}\ttfamily,
          commentstyle=\color{green}\ttfamily,
          breaklines=true,
          basicstyle=\ttfamily\footnotesize
          }
\begin{lstlisting}
//in main.cpp
#include "myclass.h"  // defines MyClass
#include "myclass.hpp"  // defines MyClass

using namespace std;

int main()
{
  std::MyClass* a;
  templates::MyClass<int>* b;
}
\end{lstlisting}
\begin{itemize}
\item Se resolve dando o nome completo!
\end{itemize}
\end{frame}

\begin{frame}[fragile]{Diretiva \textit{typedef}}
\lstset{language=C++,
          keywordstyle=\color{blue}\ttfamily,
          stringstyle=\color{red}\ttfamily,
          commentstyle=\color{green}\ttfamily,
          breaklines=true,
          basicstyle=\ttfamily\footnotesize
          }
\begin{lstlisting}
//in main.cpp
#include "myclass.h"  // defines MyClass
#include "myclass.hpp"  // defines MyClass

using namespace std;
typedef MyClass<int> MinhaClasseInteira;

int main()
{
  MyClass* a;
  MinhaClasseInteira* b;
}
\end{lstlisting}
\begin{itemize}
\item Se resolve criando uma nova definição de tipo.
\end{itemize}
\end{frame}


\begin{frame}[fragile]{Tipos Primitivos}
\lstset{language=C++,
          keywordstyle=\color{blue}\ttfamily,
          stringstyle=\color{red}\ttfamily,
          commentstyle=\color{green}\ttfamily,
          breaklines=true,
          basicstyle=\ttfamily\footnotesize
          }
\begin{itemize}
\item O tamanho é medido em word ou frações de words;
\item O tamanho é flexível e definido em cada plataforma;
\item Inteiros com e sem sinal:
\begin{itemize}
\item char, unsigned char
\item short, unsigned short
\item int, unsigned int
\item long, unsigned long
\item long long, unsigned long long
\end{itemize}
\end{itemize}
\end{frame}

\begin{frame}[fragile]{Tipos Primitivos}
\lstset{language=C++,
          keywordstyle=\color{blue}\ttfamily,
          stringstyle=\color{red}\ttfamily,
          commentstyle=\color{green}\ttfamily,
          breaklines=true,
          basicstyle=\ttfamily\footnotesize
          }
\begin{itemize}
\item Ponto flutuante:
\begin{itemize}
\item float
\item double
\item long double
\end{itemize}
\item Outros tipos:
\begin{itemize}
\item bool
\item void
\item decltype(nullptr)
\end{itemize}
\end{itemize}
\end{frame}

\begin{frame}[fragile]{Variáveis, Ponteiros e Referências}
\lstset{language=C++,
          keywordstyle=\color{blue}\ttfamily,
          stringstyle=\color{red}\ttfamily,
          commentstyle=\color{green}\ttfamily,
          breaklines=true,
          basicstyle=\ttfamily\footnotesize
          }
\begin{itemize}
\item Variável quando declarada tem um tipo que determina o seu tamanho;
\item Variável faz com que o S.O. reserve memória suficiente para o tipo e disponibilize para o programa;
\item Ponteiro é uma variável especial que guarda endereços e não dados;
\item Ponteiro quando declarado tem um tipo que determina qual o tipo do endereço que ele armazena;
\item Ponteiro faz com que o S.O. reserve memória suficiente para armazenar um endereço de memória;
\item Referência é um segundo nome para uma variável.
\end{itemize}
\end{frame}

\begin{frame}[fragile]{Variáveis, Ponteiros e Referências}
\lstset{language=C++,
          keywordstyle=\color{blue}\ttfamily,
          stringstyle=\color{red}\ttfamily,
          commentstyle=\color{green}\ttfamily,
          breaklines=true,
          basicstyle=\ttfamily\footnotesize
          }
\begin{itemize}
\item * é chamado de indireção. Trabalha com o valor do endereço de memoria guardado no ponteiro;
\item \& retorna o endereço de memória de uma variável. Também serve para criar uma refrência;
\item Podemos ter uma dupla indireção com **;
\item Podemos (em C++11) ter uma referência não constatnte com \&\&.
\end{itemize}
\end{frame}

\begin{frame}[fragile]{Variáveis, Ponteiros e Referências}
\lstset{language=C++,
          keywordstyle=\color{blue}\ttfamily,
          stringstyle=\color{red}\ttfamily,
          commentstyle=\color{green}\ttfamily,
          breaklines=true,
          basicstyle=\ttfamily\footnotesize
          }
\begin{itemize}
\item Modelo da memória de um computador:
\end{itemize}
{\rowcolors{1}{SpringGreen}{GreenYellow}
\begin{tabular}{ |p{1cm}|p{1cm}|p{1cm}|p{1cm}|p{1cm}| }
\hline
& & & & \\ \hline
& & & & \\ \hline
& & & & \\ \hline
& & & & \\ \hline
& & & & \\ \hline
\end{tabular}
}
\begin{itemize}
\item É uma sequencia de palavras de memória endereçaveis pela sua posição.
\end{itemize}
\end{frame}

\begin{frame}[fragile]{Modelo de Memória}
\lstset{language=C++,
          keywordstyle=\color{blue}\ttfamily,
          stringstyle=\color{red}\ttfamily,
          commentstyle=\color{green}\ttfamily,
          breaklines=true,
          basicstyle=\ttfamily\footnotesize
          }
\begin{itemize}
\item int a = 10;
\item Politica do S.O. Escolhe o endereço 6;
\item Modelo da memória de um computador:
\end{itemize}
{\rowcolors{1}{SpringGreen}{GreenYellow}
\begin{tabular}{ |p{1cm}|p{1cm}|p{1cm}|p{1cm}|p{1cm}| }
\hline
& & & & \\ \hline
10 & & & & \\ \hline
& & & & \\ \hline
& & & & \\ \hline
& & & & \\ \hline
\end{tabular}
}
\end{frame}

\begin{frame}[fragile]{Declarando Variáveis e Ponteiros}
\lstset{language=C++,
          keywordstyle=\color{blue}\ttfamily,
          stringstyle=\color{red}\ttfamily,
          commentstyle=\color{green}\ttfamily,
          breaklines=true,
          basicstyle=\ttfamily\footnotesize
          }
\begin{itemize}
\item int a = 10;
\item int* b = \&a;
\item Politica do S.O. Escolhe o endereço 9 para o ponteiro;
\item Modelo da memória de um computador:
\end{itemize}
{\rowcolors{1}{SpringGreen}{GreenYellow}
\begin{tabular}{ |p{1cm}|p{1cm}|p{1cm}|p{1cm}|p{1cm}| }
\hline
& & & & \\ \hline
10 & & & 6 & \\ \hline
& & & & \\ \hline
& & & & \\ \hline
& & & & \\ \hline
\end{tabular}
}
\end{frame}

\begin{frame}[fragile]{Declarando Variáveis e Ponteiros}
\lstset{language=C++,
          keywordstyle=\color{blue}\ttfamily,
          stringstyle=\color{red}\ttfamily,
          commentstyle=\color{green}\ttfamily,
          breaklines=true,
          basicstyle=\ttfamily\footnotesize
          }
\begin{lstlisting}
#include <iostream>

using namespace std;

int main(){
 int a = 5;
 int * b = &a;
 cout<<"O valor de a "<<a<<" o endereco de a e "<<&a<<endl;
 cout<<"b contem o endereco "<<b<<" o conteudo do endereco de b e: "<<*b<<endl ;
}
\end{lstlisting}
\begin{itemize}
\item Executar algumas vezes no VPL.
\end{itemize}
\end{frame}

\begin{frame}[fragile]{Declarando Variáveis e Ponteiros}
\lstset{language=C++,
          keywordstyle=\color{blue}\ttfamily,
          stringstyle=\color{red}\ttfamily,
          commentstyle=\color{green}\ttfamily,
          breaklines=true,
          basicstyle=\ttfamily\footnotesize
          }
\begin{lstlisting}
#include <iostream>

using namespace std;

int main(){
 int a = 5;
 int * b = &a;
 *b += a;
 a *= (*b);
 std::cout<<a<<std::endl;
}
\end{lstlisting}
\begin{itemize}
\item O que imprime esse código?
\end{itemize}
\end{frame}

\begin{frame}[fragile]{Referências}
\lstset{language=C++,
          keywordstyle=\color{blue}\ttfamily,
          stringstyle=\color{red}\ttfamily,
          commentstyle=\color{green}\ttfamily,
          breaklines=true,
          basicstyle=\ttfamily\footnotesize
          }
\begin{itemize}
\item Uma referência é como um "apelido" a uma outra variável;
\item O S.O. não aloca uma nova word para uma referência;
\item Tal qual um ponteiro, a referência apenas "referencia" algum endereço de memória;
\item Uma vez criada, não há como distinguir entre uma referência e o dado por ela referenciado;
\item Referência \& só aponta para um valor durante a sua vida;
\item Referência \&\& pode mudar (somente no std11);
\end{itemize}
\end{frame}

\begin{frame}[fragile]{Variáveis, Ponteiros e Referências}
\lstset{language=C++,
          keywordstyle=\color{blue}\ttfamily,
          stringstyle=\color{red}\ttfamily,
          commentstyle=\color{green}\ttfamily,
          breaklines=true,
          basicstyle=\ttfamily\footnotesize
          }
\begin{lstlisting}
int main(){
 int a = 3;
 int b = 4;
 int * p = &b;
 int & r = b;
 r += 2; // b agora vale 6
 p = &a; // p aponta para "a"
 r = a; // r ainda referencia b.
}
\end{lstlisting}
\begin{itemize}
\item Quanto vale b?
\end{itemize}
\end{frame}

\begin{frame}[fragile]{Arrays}
\lstset{language=C++,
          keywordstyle=\color{blue}\ttfamily,
          stringstyle=\color{red}\ttfamily,
          commentstyle=\color{green}\ttfamily,
          breaklines=true,
          basicstyle=\ttfamily\footnotesize
          }
\begin{lstlisting}
int main(){
 int a[5];
 a[0] = 0;
 a[4] = -3;
}
\end{lstlisting}
\begin{itemize}
\item Vetor com alocação estática;
\item O vetor inicia no endereço 7 e usa 5 variáveis inteiras;
\end{itemize}
{\rowcolors{1}{SpringGreen}{GreenYellow}
\begin{tabular}{ |p{1cm}|p{1cm}|p{1cm}|p{1cm}|p{1cm}| }
\hline
& & & & \\ \hline
& \cellcolor{Mahogany}0 &\cellcolor{Mahogany} &\cellcolor{Mahogany} & \cellcolor{Mahogany}\\ \hline
\cellcolor{Mahogany} -3 & & & & \\ \hline
& & & & \\ \hline
& & & & \\ \hline
\end{tabular}
}
\end{frame}

\begin{frame}[fragile]{Arrays}
\lstset{language=C++,
          keywordstyle=\color{blue}\ttfamily,
          stringstyle=\color{red}\ttfamily,
          commentstyle=\color{green}\ttfamily,
          breaklines=true,
          basicstyle=\ttfamily\footnotesize
          }
\begin{lstlisting}
int main(){
 int a = 5;
 int * b = new int[a];
}
\end{lstlisting}
\begin{itemize}
\item Vetor com alocação dinâmica;
\item O vetor inicia no endereço 7 e usa a variáveis inteiras;
\end{itemize}
{\rowcolors{1}{SpringGreen}{GreenYellow}
\begin{tabular}{ |p{1cm}|p{1cm}|p{1cm}|p{1cm}|p{1cm}| }
\hline
& & \cellcolor{OliveGreen}5 & & \\ \hline
& &\cellcolor{Mahogany} &\cellcolor{Mahogany} & \cellcolor{Mahogany}\\ \hline
\cellcolor{Mahogany} & \cellcolor{Mahogany} & & & \\ \hline
& & & & \\ \hline
& & & & \\ \hline
\end{tabular}
}
\end{frame}

\begin{frame}[fragile]{Iterando Arrays}
\lstset{language=C++,
          keywordstyle=\color{blue}\ttfamily,
          stringstyle=\color{red}\ttfamily,
          commentstyle=\color{green}\ttfamily,
          breaklines=true,
          basicstyle=\ttfamily\footnotesize
          }
\begin{lstlisting}
int main(){
 int a = 5;
 int * b = new int[a];
 int * it = &b[0];
 (*it++) = 8;
}
\end{lstlisting}
\begin{itemize}
\item Quais operações são executadas na última linha?\pause
\item {\color{red}(*it++) = 8 atribui o valor 8 a posição atual do iterador e incrementa o iterador, fazendo-o apontar para a próxima posição;}
\end{itemize}
\end{frame}

\begin{frame}[fragile]{Statements}
\lstset{language=C++,
          keywordstyle=\color{blue}\ttfamily,
          stringstyle=\color{red}\ttfamily,
          commentstyle=\color{green}\ttfamily,
          breaklines=true,
          basicstyle=\ttfamily\footnotesize
          }

\begin{itemize}
\item If-Then-Else:
\begin{lstlisting}
if( expressao booleana ){
 //se a expressao e verdadeira executa este bloco
} else{
 //se a expressao e falsa executa este bloco
}
\end{lstlisting}
\end{itemize}
\end{frame}

\begin{frame}[fragile]{Statements}
\lstset{language=C++,
          keywordstyle=\color{blue}\ttfamily,
          stringstyle=\color{red}\ttfamily,
          commentstyle=\color{green}\ttfamily,
          breaklines=true,
          basicstyle=\ttfamily\footnotesize
          }

\begin{itemize}
\item While:
\begin{lstlisting}
while( expressao booleana ){
 //enquanto a expressao for verdadeira executa este bloco
}
\end{lstlisting}
\item Do-While:
\begin{lstlisting}
do{
 // executa este bloco enquanto a expressao for verdadeira
}while( expressao booleana );
\end{lstlisting}
\end{itemize}
\end{frame}

\begin{frame}[fragile]{Statements}
\lstset{language=C++,
          keywordstyle=\color{blue}\ttfamily,
          stringstyle=\color{red}\ttfamily,
          commentstyle=\color{green}\ttfamily,
          breaklines=true,
          basicstyle=\ttfamily\footnotesize
          }

\begin{itemize}
\item For:
\begin{lstlisting}
for( inicializacao de variaveis ; expressao ; operacao executada a cada loop ) {
	//Para cada vez que a expressao for verdadeira, executa esse bloco. Se a expressao      for falsa, sai do bloco.
}
\end{lstlisting}
\item Em algumas versões antigas do standard tem que instanciar a variavel antes do loop.
\end{itemize}
\end{frame}

\begin{frame}[fragile]{Coleções}
\lstset{language=C++,
          keywordstyle=\color{blue}\ttfamily,
          stringstyle=\color{red}\ttfamily,
          commentstyle=\color{green}\ttfamily,
          breaklines=true,
          basicstyle=\ttfamily\footnotesize
          }

\begin{itemize}
\item Enum:
\begin{lstlisting}
enum Paises{
 Argentina,
 Brasil,
 China,
 Estados Unidos,
 Inglaterra
}
\end{lstlisting}
\begin{lstlisting}
enum Paises{
 Argentina=9,
 Brasil=6,
 China=5,
 Estados Unidos=13,
 Inglaterra=10
}
\end{lstlisting}
\end{itemize}
\end{frame}

\begin{frame}[fragile]{Switch / Case}
\lstset{language=C++,
          keywordstyle=\color{blue}\ttfamily,
          stringstyle=\color{red}\ttfamily,
          commentstyle=\color{green}\ttfamily,
          breaklines=true,
          basicstyle=\ttfamily\footnotesize
          }

\begin{itemize}
\item Seleção Switch-case:
\begin{lstlisting}
int a;
switch (a){
 case:1
  //se a vale 1, executa este bloco
  break;
 case:2
  //se a vale 2, executa este bloco
  break;
 default:
  //se a vale qualquer outro inteiro, executa este bloco
}
\end{lstlisting}
\end{itemize}
\end{frame}

{
\usebackgroundtemplate{\includegraphics[width=\paperwidth,
height=\paperheight]{../reusable_images/fundo_capa.png}}
\begin{frame}

{\LARGE Perguntas????}

\end{frame}
}


{
\usebackgroundtemplate{\includegraphics[width=\paperwidth,
height=\paperheight]{../reusable_images/fundo_capa.png}}
\begin{frame}
\includegraphics[scale=0.8]{../reusable_images/cc_logo_arge.png}\hspace{0.5cm}
\includegraphics[scale=0.95]{../reusable_images/by.png}

\vspace{1cm}
Este trabalho está licenciado sob uma Licença Creative Commons Atribuição 4.0 Internacional. Para ver uma cópia desta licença, visite http://creativecommons.org/licenses/by/4.0/.

\end{frame}
}
\end{document}
